\documentclass[aspectratio=169, 10pt]{beamer}

\usetheme{Madrid}
\usecolortheme{default}
\definecolor{ucbblue}{RGB}{0, 51, 102}

\setbeamercolor{structure}{fg=ucbblue}
\setbeamercolor*{palette primary}{fg=white,bg=ucbblue}
\setbeamercolor*{palette secondary}{fg=white,bg=ucbblue!80}
\setbeamercolor*{palette tertiary}{fg=white,bg=ucbblue!90}
\setbeamercolor{title}{fg=white, bg=ucbblue}
\setbeamercolor{frametitle}{fg=white, bg=ucbblue}
\setbeamercolor{item}{fg=ucbblue}

\usepackage[T1]{fontenc}
\usepackage[utf8]{inputenc}
\usepackage{tgpagella}
\usefonttheme{serif}
\usepackage[spanish, es-noshorthands]{babel}
\usepackage{amsmath, amssymb, bm}
\usepackage{graphicx}
\usepackage{booktabs}
\usepackage[most]{tcolorbox}
\usepackage{tikz}
\usepackage{pgfplots}
\pgfplotsset{compat=1.18}
\usepackage[american]{circuitikz}

\title[Del RC al RLC]{Unidad EC3: Del Circuito RC al RLC}
\subtitle{Modelado Físico y Matemático de Sistemas de Segundo Orden}
\author[B. Quiroga Turdera]{M.Sc. Ing. Bernardo Quiroga Turdera}
\institute[UCB Tarija]{Universidad Católica Boliviana ``San Pablo'' --- Sede Tarija}
\date{Semana 12}

\begin{document}

\begin{frame}
    \titlepage
\end{frame}

\begin{frame}{Dos Almacenamientos, Dos Variables de Estado}
    Hasta ahora estudiamos circuitos con un solo tipo de inercia (RC o RL). Si mezclamos \textbf{ambos} elementos dinámicos en una red, el sistema ahora posee "dos memorias".
    
    \vspace{0.4cm}
    \begin{columns}[T]
        \column{0.5\textwidth}
        \begin{tcolorbox}[colback=white, colframe=ucbblue, title=\textbf{El Capacitor}]
            \begin{itemize}
                \item Energía en Campo Eléctrico: $w_C = \frac{1}{2} C v_C^2$
                \item Ley Constitutiva: $i_C = C \frac{dv_C}{dt}$
                \item \textbf{Estado:} Requiere conocer $v_C(0^+)$
            \end{itemize}
        \end{tcolorbox}

        \column{0.5\textwidth}
        \begin{tcolorbox}[colback=white, colframe=ucbblue, title=\textbf{El Inductor}]
            \begin{itemize}
                \item Energía en Campo Magnético: $w_L = \frac{1}{2} L i_L^2$
                \item Ley Constitutiva: $v_L = L \frac{di_L}{dt}$
                \item \textbf{Estado:} Requiere conocer $i_L(0^+)$
            \end{itemize}
        \end{tcolorbox}
    \end{columns}
    \vspace{0.4cm}
    \textit{Para determinar la evolución completa de la red RLC, obligatoriamente necesitamos resolver \textbf{dos datos iniciales independientes}: la tensión del capacitor y la corriente del inductor.}
\end{frame}

\begin{frame}{Modelado del RLC Serie: De la Física a la Ecuación}
    \begin{columns}[T]
        \column{0.4\textwidth}
        \begin{center}
        \begin{circuitikz}[scale=0.85, transform shape]
            \draw (0,0) to[V, v=$v_s(t)$] (0,3) 
                  to[R, l=$R_{total}$, i=$i(t)$, v=$v_R$] (3,3)
                  to[L, l=$L$, v=$v_L$] (3,0)
                  to[C, l=$C$, v=$v_C$] (0,0);
        \end{circuitikz}
        \end{center}
        
        \column{0.6\textwidth}
        1. Aplicamos Ley de Voltajes de Kirchhoff (KVL):
        $$v_s(t) = v_R(t) + v_L(t) + v_C(t)$$
        
        2. Sustituimos Ley de Ohm y Ley del Inductor:
        $$v_s = R_{total} \cdot i + L \frac{di}{dt} + v_C$$
        
        3. Identificamos la corriente del lazo por la Ley del Capacitor ($i = i_C = C \frac{dv_C}{dt}$). Por tanto, la derivada de la corriente es:
        $$\frac{di}{dt} = C \frac{d^2v_C}{dt^2}$$
        
        4. Reemplazamos en KVL para unificar variables:
        $$v_s = R_{total} \left( C \frac{dv_C}{dt} \right) + L \left( C \frac{d^2v_C}{dt^2} \right) + v_C$$
    \end{columns}
\end{frame}

\begin{frame}{La Ecuación Diferencial de Segundo Orden}
    Reordenando y dividiendo toda la expresión entre $LC$ (asumiendo $L>0, C>0$), obtenemos la forma canónica de la EDO de un circuito RLC Serie:
    
    \begin{alertblock}{Modelo Matemático RLC Serie}
        $$ \frac{d^2v_C}{dt^2} + \left( \frac{R_{total}}{L} \right) \frac{dv_C}{dt} + \left( \frac{1}{LC} \right) v_C = \frac{1}{LC} v_s(t) $$
    \end{alertblock}
    
    \vspace{0.3cm}
    \textbf{Interpretación Física de los Términos:}
    \begin{itemize}
        \item Término de $\frac{d^2v_C}{dt^2}$: Inercia dinámica asociada al inductor (almacenamiento magnético).
        \item Término de $\frac{dv_C}{dt}$ (con $R_{total}$): \textbf{Disipación} (Amortiguamiento de la energía en forma de calor por la resistencia de lazo).
        \item Término $v_C$: Estado elástico de restitución asociado al capacitor.
    \end{itemize}
\end{frame}

\begin{frame}{Parámetros de Diseño: Frecuencia Natural y Amortiguamiento}
    Para simplificar la ecuación, definimos dos parámetros fundamentales de control:
    
    \begin{columns}[T]
        \column{0.5\textwidth}
        \begin{tcolorbox}[colback=white, colframe=black, title=\textbf{Frecuencia Natural no Amortiguada}]
            $$\omega_0 = \frac{1}{\sqrt{LC}} \quad [\text{rad/s}]$$
            \textit{¿A qué velocidad intercambiarían energía $L$ y $C$ si la resistencia del circuito fuera literalmente cero (sin pérdida de calor)?} \\
            $f_0 = \frac{\omega_0}{2\pi} \quad [\text{Hz}]$
        \end{tcolorbox}
        
        \column{0.5\textwidth}
        \begin{tcolorbox}[colback=white, colframe=black, title=\textbf{Factor de Amortiguamiento}]
            $$\alpha = \frac{R_{total}}{2L} \quad [\text{s}^{-1}]$$
            \textit{¿Qué tan rápido se disipa o "frena" la energía en la resistencia?}
        \end{tcolorbox}
    \end{columns}
    
    \vspace{0.4cm}
    La EDO puede reescribirse elegantemente como:
    $$ v_C'' + 2\alpha v_C' + \omega_0^2 v_C = \omega_0^2 v_s(t) $$
\end{frame}

\begin{frame}{Regímenes de Amortiguamiento (Factor $\zeta = \alpha / \omega_0$)}
    La relación entre pérdida resistiva ($\alpha$) y energía reactiva ($\omega_0$) define la forma de la respuesta:
    
    \vspace{0.2cm}
    \begin{itemize}
        \item \textbf{Subamortiguado ($\zeta < 1$):} Poca resistencia. La energía "rebota" entre el capacitor y el inductor causando \textbf{sobreimpulso} y \textbf{oscilaciones} decrecientes. 
        \item \textbf{Crítico ($\zeta = 1$):} Ocurre en $R_{crit} = 2\sqrt{\frac{L}{C}}$. Es la resistencia de transición donde el sistema retorna al estado final sin oscilar.
        \item \textbf{Sobreamortiguado ($\zeta > 1$):} Exceso de resistencia. El circuito es sumamente viscoso, no oscila, y llega asintóticamente a la meta (usualmente más lento que el crítico).
    \end{itemize}
    
    \vspace{0.3cm}
    \textit{Importante: ¡Tener inductores y capacitores juntos NO garantiza oscilación! Si $R_{total}$ es alta, el sistema aniquilará la oscilación antes de que empiece.}
\end{frame}

\begin{frame}{Clasificación Nominal (Ejemplo de Aplicación)}
    Considere un circuito con $L = 10\,\text{mH}$ y $C = 1\,\mu\text{F}$. 
    \begin{itemize}
        \item Frecuencia natural: $\omega_0 = \frac{1}{\sqrt{0.01 \cdot 10^{-6}}} = 10,000\,\text{rad/s}$ ($f_0 \approx 1.59\,\text{kHz}$).
        \item Resistencia crítica total: $R_{crit} = 2\sqrt{\frac{10\text{mH}}{1\mu\text{F}}} = \mathbf{200\,\Omega}$.
    \end{itemize}

    \vspace{0.2cm}
    \renewcommand{\arraystretch}{1.3}
    \begin{tabular}{cccc}
    \toprule
    \textbf{$R_{total}$ hipotética} & $\alpha = \frac{R_{total}}{2L}$ & $\zeta = \frac{R_{total}}{R_{crit}}$ & \textbf{Régimen} \\
    \midrule
    $50\,\Omega$ & $2500\,\text{s}^{-1}$ & $0.25$ & \textbf{Subamortiguado (Oscila)} \\
    $200\,\Omega$ & $10000\,\text{s}^{-1}$ & $1.00$ & \textbf{Críticamente Amortiguado} \\
    $400\,\Omega$ & $20000\,\text{s}^{-1}$ & $2.00$ & \textbf{Sobreamortiguado (Lento)} \\
    \bottomrule
    \end{tabular}
    
    \vspace{0.3cm}
    \begin{tcolorbox}[colback=gray!10, colframe=black, title=Física del Mundo Real]
    ¡Atención! $R_{total}$ es la resistencia total del lazo, no solo el resistor externo comprado. \\
    $R_{total} = R_{externa} + R_{generador} + R_{DCR\_bobina} + R_{cables}$.
    \end{tcolorbox}
\end{frame}

\begin{frame}{Formas de Onda (Comparativa Visual Teórica)}
    \begin{center}
    \begin{tikzpicture}
        \begin{axis}[
            width=12cm, height=6.5cm,
            xlabel={Tiempo $t$ (ms)}, ylabel={$v_C(t)$ (V)},
            xmin=0, xmax=1.5, ymin=0, ymax=15,
            grid=major,
            legend pos=north east,
            legend style={font=\small}
        ]
        % Curvas conceptuales (v_C(infty) = 10V)
        % Subamortiguado (zeta=0.25, w0=10000)
        \addplot[blue, thick, domain=0:1.5, samples=200] {10 - 10*exp(-2.5*x)* (cos(deg(9.68*x)) + (2.5/9.68)*sin(deg(9.68*x)))};
        \addlegendentry{Subamortiguado ($50\Omega$)}
        
        % Critico (zeta=1, w0=10000)
        \addplot[black, thick, domain=0:1.5, samples=100] {10 - 10*(1 + 10*x)*exp(-10*x)};
        \addlegendentry{Crítico ($200\Omega$)}
        
        % Sobreamortiguado (zeta=2, w0=10000) -> s1=-2.68, s2=-37.32
        \addplot[red, thick, domain=0:1.5, samples=100] {10 - 10*(1.077*exp(-2.68*x) - 0.077*exp(-37.32*x))};
        \addlegendentry{Sobreamortiguado ($400\Omega$)}
        
        \addplot[dashed, gray] coordinates {(0,10) (1.5,10)};
        \end{axis}
    \end{tikzpicture}
    \end{center}
\end{frame}

\end{document}
